\chapter{Tobramycin Level}

\section*{What Is This Test?}
The tobramycin level measures the serum concentration of tobramycin, an aminoglycoside antibiotic closely related to gentamicin with enhanced activity against Pseudomonas aeruginosa (generally 2--4 times more potent than gentamicin against P. aeruginosa). Tobramycin is used for serious Gram-negative infections, particularly in patients with cystic fibrosis (CF) who have chronic P. aeruginosa respiratory infections, severe nosocomial infections, febrile neutropenia with suspected Gram-negative sepsis, and in combination with beta-lactam antibiotics for synergistic treatment of endocarditis. Like gentamicin, tobramycin binds to the 30S ribosomal subunit, inhibiting protein synthesis. Tobramycin is available for IV administration and as an inhaled formulation (tobramycin inhalation solution [TOBI] and tobramycin inhalation powder [TOBI Podhaler]) for chronic suppressive therapy in CF patients with chronic P. aeruginosa infection. Therapeutic drug monitoring (TDM) is essential for systemic (IV) tobramycin due to its narrow therapeutic index: nephrotoxicity (non-oliguric acute tubular necrosis, 10--25\%), ototoxicity (irreversible vestibular and cochlear toxicity, 5--15\%), and neuromuscular blockade (rare, in myasthenia gravis or with concurrent neuromuscular blocking agents). Extended-interval (once-daily) dosing is preferred for most indications except for synergy in enterococcal endocarditis.

\section*{Alternative Names}
Tobramycin Level; Tobramycin Peak and Trough; Nebcin Level; TOBI Level; Aminoglycoside Level; Tobramycin Serum Concentration

\section*{Principle}
Tobramycin is measured by particle-enhanced turbidimetric inhibition immunoassay (PETINIA) or enzyme-multiplied immunoassay technique (EMIT) using specific anti-tobramycin antibodies. Tobramycin in patient serum competes with a labeled tobramycin-conjugated particle for antibody binding. The rate of aggregation is inversely proportional to the tobramycin concentration.

\section*{History}
Tobramycin was discovered in the 1960s from the actinomycete Streptoalloteichus tenebrarius (formerly Streptomyces tenebrarius). It received FDA approval in 1975. Inhaled tobramycin for cystic fibrosis was developed in the 1990s and approved in 1997 (TOBI), representing a major advance in the chronic suppressive management of P. aeruginosa in CF, reducing the frequency of pulmonary exacerbations and improving lung function (FEV1).

\section*{Why This Test Is Ordered}
TDM for IV tobramycin to maximize efficacy and minimize toxicity. Especially important in: CF patients (altered pharmacokinetics: larger volume of distribution, increased renal clearance requiring higher doses per kg body weight), critically ill patients, renal impairment, prolonged therapy, burns, and concurrent nephrotoxic drugs.

\section*{Primary vs Secondary Test}
This is a primary therapeutic drug monitoring test.

\section*{How This Test Is Performed}
Blood sample in SST (gold-top) or red-top tube. Peak: 30 minutes after IV infusion end. Trough: immediately before next dose. For extended-interval dosing: single level 6--14 hours after dose, plotted on Hartford nomogram.

\section*{What Preparation Is Needed}
Accurate documentation of the dose and the time of the blood draw is essential. No fasting is required. The sample should be drawn at the correct time relative to the dose: the peak 30 minutes after the end of the IV infusion, the trough immediately before the next dose, or, for extended-interval (once-daily) dosing, a single level 6--14 hours after the dose. Levels should be obtained after steady state has been reached (approximately after 2--3 doses). The exact times of the last dose and the blood draw must be recorded on the requisition.

\section*{Contraindications and Precautions}
There are no contraindications to the blood draw. Relative contraindications to tobramycin therapy itself: a history of severe hypersensitivity to an aminoglycoside, myasthenia gravis or other neuromuscular disease (increased risk of neuromuscular blockade), and pregnancy (possible fetal eighth-nerve toxicity; use only when clearly indicated). Precautions: The Hartford nomogram should NOT be used in patients with cystic fibrosis, pregnancy, burns or ascites, or creatinine clearance <20 mL/min, because their volume of distribution and clearance differ from the nomogram's reference population. Concurrent nephrotoxic drugs (amphotericin B, vancomycin, loop diuretics, radiocontrast) increase the risk of nephrotoxicity, and a baseline serum creatinine should be obtained before the first dose.

\section*{Risks}
Minimal risk from the venipuncture: slight pain or bruising at the site, and rarely hematoma or vasovagal syncope. The significant toxicities of tobramycin (nephrotoxicity and irreversible ototoxicity) are caused by the drug itself, not the blood draw, and are precisely why the level is monitored.

\section*{What Happens After the Test}
Results are typically available within hours and are interpreted in relation to the dosing regimen, the indication, and renal function. If the trough is elevated or the peak is subtherapeutic, the dose and/or dosing interval is adjusted, and a repeat level is drawn after a new steady state is reached (after 2--3 doses). Serum creatinine and electrolytes are monitored at least 2--3 times weekly, and audiometry is arranged for prolonged courses or high-risk patients.

\section*{Understanding Results}

\subsection*{Subtherapeutic}
Peak <4--5 mcg/mL (multiple-daily dosing). Inadequate bactericidal effect and risk of treatment failure.

\subsection*{Potentially Toxic}
Trough >2 mcg/mL (multiple-daily dosing) or >0.5--1 mcg/mL (extended-interval dosing). Increased risk of nephrotoxicity (acute tubular necrosis presenting as non-oliguric rise in serum creatinine after 5--7 days, generally reversible) and ototoxicity (irreversible; affects high-frequency hearing first, progressing to conversational frequencies; vestibular damage causes oscillopsia and ataxia).

\section*{Normal Results}
Traditional multiple-daily dosing: Peak 4--10 mcg/mL (severe Gram-negative infections), 3--5 mcg/mL (synergy for endocarditis). Trough <2 mcg/mL (<1 mcg/mL preferred). Extended-interval dosing: Peak 15--25 mcg/mL. Trough <0.5--1 mcg/mL. Cystic fibrosis patients: higher doses often required (10--15 mg/kg/day divided or once-daily). The Hartford nomogram should NOT be used in CF patients, pregnancy, burns/ascites, or CrCl <20 mL/min.

\section*{Follow-up Tests}
Serum creatinine/BUN (baseline and at least 2--3 times weekly), serum electrolytes (magnesium, potassium), audiometry for prolonged courses or high-risk patients.

\section*{References}
\begin{enumerate}
  \item Avent ML, Rogers BA, Cheng AC, Paterson DL. Current use of aminoglycosides. Intern Med J. 2011;41(6):441--449.
  \item Ramsey BW, Pepe MS, Quan JM, et al. Intermittent administration of inhaled tobramycin in patients with cystic fibrosis. N Engl J Med. 1999;340(1):23--30.
  \item Nicolau DP, Freeman CD, Belliveau PP, et al. Experience with a once-daily aminoglycoside program in 2,184 adult patients. Antimicrob Agents Chemother. 1995;39(3):650--655.
  \item Winter ME. Basic Clinical Pharmacokinetics. 5th ed. Baltimore: Lippincott Williams \& Wilkins; 2010.
  \item Burton ME, Shaw LM, Schentag JJ, Evans WE. Applied Pharmacokinetics \& Pharmacodynamics: Principles of Therapeutic Drug Monitoring. 4th ed. Baltimore: Lippincott Williams \& Wilkins; 2006.
\end{enumerate}

\clearpage
\section*{Regulatory Status and Authorization Date}
Regulatory status applies to the specific assay, instrument, manufacturer, and intended use rather than to the test name alone. No single approval date can be assigned to this test category. For a commercial assay, verify the current product in the relevant regulator's database: the U.S. Food and Drug Administration (FDA) clearance, approval, or authorization; India's Central Drugs Standard Control Organisation (CDSCO) licence or permission; the European Union In Vitro Diagnostic Regulation (EU IVDR) CE certificate issued through the applicable conformity-assessment route; or the United Kingdom Medicines and Healthcare products Regulatory Agency (MHRA) registration and UKCA/CE status. Laboratory-developed tests may not have an individual FDA, CDSCO, EU IVDR, or MHRA approval date.
\clearpage
